\UseRawInputEncoding
\documentclass[12pt,a4paper]{article}

% Packages for formatting and graphics
\usepackage{amsmath}
\usepackage{graphicx}
\usepackage{hyperref}
\usepackage{float}
\hypersetup{
    colorlinks=true,
    linkcolor=black,
    urlcolor=blue,
    citecolor=blue
}
\usepackage{geometry}
\geometry{a4paper, margin=1in}

\usepackage{listings}
\usepackage{color}
\usepackage{ragged2e}

\definecolor{codegreen}{rgb}{0,0.6,0}
\definecolor{codegray}{rgb}{0.5,0.5,0.5}
\definecolor{codepurple}{rgb}{0.58,0,0.82}
\definecolor{backcolour}{rgb}{0.95,0.95,0.92}

\lstdefinestyle{mystyle}{
    backgroundcolor=\color{backcolour},   
    commentstyle=\color{codegreen},
    keywordstyle=\color{magenta},
    numberstyle=\tiny\color{codegray},
    stringstyle=\color{codepurple},
    basicstyle=\ttfamily\footnotesize,
    breakatwhitespace=false,         
    breaklines=true,                 
    captionpos=b,                    
    keepspaces=true,                 
    numbers=left,                    
    numbersep=5pt,                  
    showspaces=false,                
    showstringspaces=false,
    showtabs=false,                  
    tabsize=2
}


% Title page
\title{EN3150 Assignment 01: Learning from Data and Related Challenges and Linear Models for Regression}
\author{Your Name \\ Your Index Number}
\date{\today}

\begin{document}

\begin{titlepage}
    \begin{center}
        \vspace*{0.01cm}
        \LARGE
        \textbf{Department of Electronic and Telecommunication Engineering\\
        University of Moratuwa}
        
        \vspace{1.2cm}
        \LARGE
        \textbf{EN3150 - Pattern Recognition}
        
        \vspace{1.0cm}

        \includegraphics[width=0.4\textwidth]{university.jpeg}

        \vspace{1.5cm}
        \LARGE
        Assignment 03 \\
        Simple convolutional neural network to perform classification
        \vspace{5cm}
        
        \begin{tabular}{cc}
            \Large 230415H & \Large Mudaliarachchi N.S. \\
            \Large 230185B    & \Large FERNANDO S.D.     \\
        \end{tabular}

    \end{center}
    \begin{center}
    \small 05/09/2026
    \end{center}
\end{titlepage}
\newpage

% Table of Contents
\tableofcontents
\newpage

\listoftables
\newpage

% Section 1: 
\section{ Resource-Constrained CNN for Edge Image Classification }

\subsection{ModelA(Standard CNN)}

\begin{lstlisting}[language=python, style=mystyle]
def hello():
    return "hi"
print(hello())
\end{lstlisting}


\subsection{ ModelB(Lightweight CNN)}
\subsection{Optimizer Selection & Tuning}
\subsection{Custom Model Training & Evaluation}
\newpage
\section{ Compare with Lightweight State-of-the-Art Networks}

\subsection{SOTAFine-Tuning & Evaluation}

Two pre trained lightweight convolutional neural network architectures, MobileNetV2 and EfficientNet-B0, were selected for evaluation. Both architectures are designed to provide efficient image classification and are suitable for resource constrained and mobile/edge computing environments. The same training, validation, and test splits were used for both models to ensure a consistent comparison. The dataset was divided into 70\% training, 15\% validation, and 15\% testing subsets using the same stratified split. Both models used ImageNet pretrained weights, and input images were resized to $224\times224$ pixels.

\subsubsection{MobileNetV2}

MobileNetV2 was first loaded with its ImageNet pretrained weights. Since the original model was trained to classify 1000 ImageNet classes, its final classification layer was replaced with a new fully connected layer with 9 output neurons, corresponding to the nine waste categories in the RealWaste dataset.

Initially, the feature extraction layers of MobileNetV2 were frozen, while only the newly added classification layer was trained. The model was trained for 10 epochs using the Adam optimizer with a learning rate of $0.001$. This allowed the newly added classifier to adapt the pretrained ImageNet features to the waste classification task without modifying the pretrained feature extractor.

After the initial 10 epochs, the deeper layers of the network were fine-tuned. Specifically, feature blocks 14--18 were unfrozen while the earlier feature layers remained frozen. The learning rate was reduced to $0.0001$ to allow the pretrained weights to be adjusted gradually without significantly disrupting the previously learned features. The model was subsequently fine tuned for a further 5 epochs.

The best validation accuracy obtained during fine tuning was $89.20\%$. The final model was then evaluated on the previously unseen test set, obtaining a test loss of $0.3444$ and a test accuracy of $88.50\%$.

\begin{table}[H]
    \centering
    \caption{MobileNetV2 classification report on the RealWaste test set.}
    \label{tab:mobilenetv2_classification}
    \begin{tabular}{lcccc}
        \hline
        Class & Precision & Recall & F1-score & Support \\
        \hline
        Cardboard & 0.94 & 0.93 & 0.93 & 69 \\
        Food Organics & 0.93 & 0.92 & 0.93 & 62 \\
        Glass & 0.85 & 0.89 & 0.87 & 63 \\
        Metal & 0.85 & 0.89 & 0.87 & 118 \\
        Miscellaneous Trash & 0.77 & 0.79 & 0.78 & 75 \\
        Paper & 0.99 & 0.88 & 0.93 & 75 \\
        Plastic & 0.87 & 0.84 & 0.86 & 138 \\
        Textile Trash & 0.88 & 0.96 & 0.92 & 48 \\
        Vegetation & 0.95 & 0.95 & 0.95 & 65 \\
        \hline
        Accuracy & & & 0.88 & 713 \\
        Macro Avg. & 0.89 & 0.89 & 0.89 & 713 \\
        Weighted Avg. & 0.89 & 0.88 & 0.89 & 713 \\
        \hline
    \end{tabular}
\end{table}

The confusion matrix for MobileNetV2 is shown in Table~\ref{tab:mobilenetv2_confusion}.

\begin{table}[H]
    \centering
    \caption{MobileNetV2 confusion matrix on the RealWaste test set.}
    \label{tab:mobilenetv2_confusion}
    \resizebox{\textwidth}{!}{
    \begin{tabular}{lccccccccc}
        \hline
        & Cardboard & Food Org. & Glass & Metal & Misc. Trash & Paper & Plastic & Textile & Vegetation \\
        \hline
        Cardboard
        & 64 & 0 & 2 & 2 & 0 & 1 & 0 & 0 & 0 \\
        
        Food Organics
        & 0 & 57 & 0 & 0 & 1 & 0 & 2 & 0 & 2 \\
        
        Glass
        & 1 & 0 & 56 & 3 & 2 & 0 & 1 & 0 & 0 \\
        
        Metal
        & 0 & 0 & 1 & 105 & 5 & 0 & 7 & 0 & 0 \\
        
        Misc. Trash
        & 0 & 2 & 1 & 1 & 59 & 0 & 6 & 5 & 1 \\
        
        Paper
        & 2 & 1 & 1 & 2 & 2 & 66 & 1 & 0 & 0 \\
        
        Plastic
        & 1 & 1 & 5 & 11 & 3 & 0 & 116 & 1 & 0 \\
        
        Textile Trash
        & 0 & 0 & 0 & 0 & 2 & 0 & 0 & 46 & 0 \\
        
        Vegetation
        & 0 & 0 & 0 & 0 & 3 & 0 & 0 & 0 & 62 \\
        \hline
    \end{tabular}
    }
\end{table}

The model correctly classified most samples in the Cardboard, Food Organics, Paper, Textile Trash, and Vegetation categories. The most noticeable confusion occurred between Plastic and Metal, where 11 Plastic samples were classified as Metal and 7 Metal samples were classified as Plastic. Miscellaneous Trash also showed relatively greater confusion with Plastic and Textile Trash.

\subsubsection{EfficientNet-B0}

EfficientNet-B0 was subsequently evaluated using the same dataset split and preprocessing procedure. The model was initialized using ImageNet-pretrained weights. Similar to MobileNetV2, the original 1000-class classification layer was replaced with a new fully connected layer producing 9 output classes.

Initially, the EfficientNet-B0 feature extractor was frozen and only the newly added classifier was trained. The model was trained for 10 epochs using the Adam optimizer with a learning rate of $0.001$. At the end of this stage, the model achieved a validation accuracy of $83.03\%$.

The deeper feature layers were then fine-tuned by unfreezing feature blocks 6--8 while keeping the earlier feature blocks frozen. The learning rate was reduced to $0.0001$, and the model was fine tuned for an additional 5 epochs. Validation accuracy increased to $89.90\%$ at the end of this stage.

To investigate whether additional fine tuning could improve the model, a further 5 epochs were performed using the same learning rate and trainable layers. The resulting model achieved a test loss of $0.3521$ and a test accuracy of $88.92\%$.

Thus, the EfficientNet-B0 experiment followed the same two stage training strategy used for MobileNetV2: an initial classifier training stage with the pretrained backbone frozen, followed by low learning rate fine tuning of the deeper feature layers.


\begin{table}[H]
    \centering
    \caption{EfficientNet-B0 classification report on the RealWaste test set.}
    \label{tab:efficientnet_classification}
    \begin{tabular}{lcccc}
        \hline
        Class & Precision & Recall & F1-score & Support \\
        \hline
        Cardboard & 0.92 & 0.94 & 0.93 & 69 \\
        Food Organics & 0.93 & 0.90 & 0.92 & 62 \\
        Glass & 0.93 & 0.84 & 0.88 & 63 \\
        Metal & 0.87 & 0.89 & 0.88 & 118 \\
        Miscellaneous Trash & 0.79 & 0.80 & 0.79 & 75 \\
        Paper & 0.94 & 0.88 & 0.91 & 75 \\
        Plastic & 0.85 & 0.89 & 0.87 & 138 \\
        Textile Trash & 0.90 & 0.92 & 0.91 & 48 \\
        Vegetation & 0.97 & 0.95 & 0.96 & 65 \\
        \hline
        Accuracy & & & 0.89 & 713 \\
        Macro Avg. & 0.90 & 0.89 & 0.89 & 713 \\
        Weighted Avg. & 0.89 & 0.89 & 0.89 & 713 \\
        \hline
    \end{tabular}
\end{table}

The confusion matrix for EfficientNet-B0 is shown in Table~\ref{tab:efficientnet_confusion}. 

\begin{table}[H]
    \centering
    \caption{EfficientNet-B0 confusion matrix on the RealWaste test set.}
    \label{tab:efficientnet_confusion}
    \resizebox{\textwidth}{!}{
    \begin{tabular}{lccccccccc}
        \hline
        & Cardboard & Food Org. & Glass & Metal & Misc. Trash & Paper & Plastic & Textile & Vegetation \\
        \hline
        Cardboard
        & 65 & 0 & 0 & 0 & 0 & 1 & 3 & 0 & 0 \\
        
        Food Organics
        & 0 & 56 & 0 & 0 & 4 & 0 & 0 & 1 & 1 \\
        
        Glass
        & 0 & 0 & 53 & 3 & 2 & 1 & 4 & 0 & 0 \\
        
        Metal
        & 1 & 0 & 0 & 105 & 4 & 1 & 7 & 0 & 0 \\
        
        Misc. Trash
        & 1 & 1 & 1 & 3 & 60 & 1 & 5 & 3 & 0 \\
        
        Paper
        & 4 & 1 & 0 & 1 & 0 & 66 & 2 & 1 & 0 \\
        
        Plastic
        & 0 & 0 & 3 & 8 & 3 & 0 & 123 & 0 & 1 \\
        
        Textile Trash
        & 0 & 0 & 0 & 1 & 2 & 0 & 1 & 44 & 0 \\
        
        Vegetation
        & 0 & 2 & 0 & 0 & 1 & 0 & 0 & 0 & 62 \\
        \hline
    \end{tabular}
    }
\end{table}

The EfficientNet-B0 model correctly classified most samples across the nine waste categories. The strongest classification performance was observed for Vegetation, with an F1-score of $0.96$, followed by Cardboard with an F1-score of $0.93$. The largest number of misclassifications occurred between Plastic and Metal, where 8 Plastic samples were classified as Metal and 7 Metal samples were classified as Plastic. Miscellaneous Trash also showed some confusion with Plastic and Textile Trash.



\begin{table}[H]
    \centering
    \caption{Test performance comparison of the pre-trained lightweight architectures.}
    \label{tab:pretrained_performance}
    \begin{tabular}{lcc}
        \hline
        Model & Test Loss & Test Accuracy \\
        \hline
        MobileNetV2 & 0.3444 & 88.50\% \\
        EfficientNet-B0 & 0.3521 & 88.92\% \\
        \hline
    \end{tabular}
\end{table}

\begin{table}[H]
    \centering
    \caption{Model complexity of the pre-trained lightweight architectures.}
    \label{tab:model_complexity}
    \begin{tabular}{lcc}
        \hline
        Model & Total Parameters & Saved Model Size (MB) \\
        \hline
        MobileNetV2 & 2,235,401 & 8.76 \\
        EfficientNet-B0 & 4,019,077 & 15.62 \\
        \hline
    \end{tabular}
\end{table}
\newpage
\subsection{Final Comparison & Trade-off Discussion}

% End of document
\end{document}
